\begin{figure}[H]
    \centering

    %--------------------- DIVING ---------------------%
    \begin{subfigure}{0.32\textwidth}
        \centering
        \begin{tikzpicture}[
            scale=1.0,
            base/.style={circle,draw=gray!60,fill=white,inner sep=1pt,minimum size=5mm},
            treeedge/.style={draw=gray!40},
            feasible/.style={fill=green!30,draw=green!60!black},
            infeasible/.style={fill=red!20,draw=red!70!black,dashed},
            pathDiving/.style={very thick,->,>=Stealth,blue},
        ]
            %--- Nodes (same layout for all subfigures) ---%
            \node[base,infeasible] (r)  at (0,0) {};
            \node[base,infeasible] (a)  at (-1.5,-1.5) {};
            \node[base,infeasible] (b)  at ( 1.5,-1.5) {};
            \node[base,infeasible] (a1) at (-2.3,-3) {};
            \node[base,feasible] (a2) at (-0.7,-3) {};
            \node[base,feasible] (b1) at ( 0.7,-3) {};
            \node[base,infeasible] (b2) at ( 2.3,-3) {};

            %--- Tree edges (faint) ---%
            \draw[treeedge] (r) -- (a);
            \draw[treeedge] (r) -- (b);
            \draw[treeedge] (a) -- (a1);
            \draw[treeedge] (a) -- (a2);
            \draw[treeedge] (b) -- (b1);
            \draw[treeedge] (b) -- (b2);

            %--- Diving path: root -> left child -> left grandchild ---%
            \draw[pathDiving] (r) to[bend left=5] (a);
            \draw[pathDiving] (a) to[bend left=5] (a1);

        \end{tikzpicture}
        \caption{Diving: single guided root-to-leaf path}
        \label{fig:diving_path}
    \end{subfigure}
    \hfill
    %--------------------- DFS ---------------------%
    \begin{subfigure}{0.32\textwidth}
        \centering
        \begin{tikzpicture}[
            scale=1.0,
            base/.style={circle,draw=gray!60,fill=white,inner sep=1pt,minimum size=5mm},
            feasible/.style={fill=green!30,draw=green!60!black},
            infeasible/.style={fill=red!20,draw=red!70!black,dashed},
            treeedge/.style={draw=gray!40},
            pathDFS/.style={very thick,->,>=Stealth,red},
        ]
            %--- Nodes ---%
            \node[base,infeasible] (r)  at (0,0) {};
            \node[base,infeasible] (a)  at (-1.5,-1.5) {};
            \node[base,infeasible] (b)  at ( 1.5,-1.5) {};
            \node[base,infeasible] (a1) at (-2.3,-3) {};
            \node[base,feasible] (a2) at (-0.7,-3) {};
            \node[base,feasible] (b1) at ( 0.7,-3) {};
            \node[base,infeasible] (b2) at ( 2.3,-3) {};

            %--- Tree edges (faint) ---%
            \draw[treeedge] (r) -- (a);
            \draw[treeedge] (r) -- (b);
            \draw[treeedge] (a) -- (a1);
            \draw[treeedge] (a) -- (a2);
            \draw[treeedge] (b) -- (b1);
            \draw[treeedge] (b) -- (b2);

            %--- DFS path (with backtracking, using bent arrows) ---%
            % r -> a -> a1
            \draw[pathDFS] (r)  to[bend left=5] (a);
            \draw[pathDFS] (a)  to[bend left=5] (a1);
            % backtrack a1 -> a
            %\draw[pathDFS] (a1) to[bend left=40] (a);
            % a -> a2
            \draw[pathDFS] (a1)  to[bend left=5] (a2);


        \end{tikzpicture}
        \caption{Depth-first search with backtracking}
        \label{fig:dfs_path}
    \end{subfigure}
    \hfill
    %--------------------- BFS ---------------------%
    \begin{subfigure}{0.32\textwidth}
        \centering
        \begin{tikzpicture}[
            scale=1.0,
            base/.style={circle,draw=gray!60,fill=white,inner sep=1pt,minimum size=5mm},
            feasible/.style={fill=green!30,draw=green!60!black},
            infeasible/.style={fill=red!20,draw=red!70!black,dashed},
            treeedge/.style={draw=gray!40},
            pathBFS/.style={very thick,->,>=Stealth,green!60!black},
        ]
            %--- Nodes ---%
            \node[base,infeasible] (r)  at (0,0) {};
            \node[base,infeasible] (a)  at (-1.5,-1.5) {};
            \node[base,infeasible] (b)  at ( 1.5,-1.5) {};
            \node[base,infeasible] (a1) at (-2.3,-3) {};
            \node[base,feasible] (a2) at (-0.7,-3) {};
            \node[base,feasible] (b1) at ( 0.7,-3) {};
            \node[base,infeasible] (b2) at ( 2.3,-3) {};

            %--- Tree edges (faint) ---%
            \draw[treeedge] (r) -- (a);
            \draw[treeedge] (r) -- (b);
            \draw[treeedge] (a) -- (a1);
            \draw[treeedge] (a) -- (a2);
            \draw[treeedge] (b) -- (b1);
            \draw[treeedge] (b) -- (b2);

            %--- BFS path (level by level, with bent backtracking to root/parent) ---%
            % Level 0 -> Level 1
            %\draw[pathBFS] (r)  to[bend left=5] (a);
            %\draw[pathBFS] (a)  to[bend left=40] (r);
            %\draw[pathBFS] (r)  to[bend left=5] (b);
            %\draw[pathBFS] (b)  to[bend left=40] (r);
            % Level 1 -> Level 2
            %\draw[pathBFS] (r)  to[bend left=20] (a1);
            %\draw[pathBFS] (a1) to[bend left=40] (r);
            %\draw[pathBFS] (r)  to[bend left=0]  (a2);

            \draw[pathBFS] (r) to[bend left=5] (a);
            \draw[pathBFS] (a) to (b);
            \draw[pathBFS] (b) to (a1);
            \draw[pathBFS] (a1) to (a2);

            

        \end{tikzpicture}
        \caption{Breadth-first search visiting nodes by level}
        \label{fig:bfs_path}
    \end{subfigure}

    \caption{Schematic comparison of Diving (blue), depth-first search (red), and breadth-first search (green) traversal patterns on the same search tree. Faint gray edges show the full tree; colored arrows indicate the order in which nodes are visited.}
    \label{fig:diving_dfs_bfs}
\end{figure}